\documentclass{article}
\usepackage{amsmath, amssymb, amsthm}
\usepackage{hyperref}
\usepackage{geometry}
\geometry{margin=1in}

\title{Erd\H{o}s Problem \#809}
\date{Last updated: 2025-08-31}
\author{Source: erdosproblems.com}

\begin{document}
\maketitle

\noindent\textbf{Status:} OPEN \quad \textbf{No prize}

\noindent\textbf{Tags:} graph theory, ramsey theory

\medskip
\noindent\textbf{Problem Statement:}

Define the anti-Ramsey number $\chi_S(n,e,G)$ as the smallest $r$ such that there is a graph with $n$ vertices and $e$ edges with an $r$-colouring of its edges in which every copy of $G$ has entirely distinct edge colours.

Is it true that, for all $k\geq 3$,\[\chi_S(n, \lfloor n^2/4\rfloor+1,C_{2k+1})\sim n^2/8?\]

\bigskip
\noindent\textbf{Remarks:}

A problem of Burr, Erd\H{o}s, Graham, and S\'{o}s \cite{BEGS89} who proved that\[\chi_S(n, \lfloor n^2/4\rfloor+1,C_{2k+1})\gg_k n^2.\]This question was solved in the affirmative for all $k\geq 4$ by Buci\'{c}, Chen, and Ma \cite{BCM26}.

The situation for $C_3$ and $C_5$ are quite different - it is easy to see that\[\chi_S(n, \lfloor n^2/4\rfloor+1,C_{3})=3\]and Erd\H{o}s and Simonovits proved (as reported in \cite{BEGS89}) that\[\chi_S(n, \lfloor n^2/4\rfloor+1,C_{5})=\lfloor n/2\rfloor+3\]for all large $n$.

See also [810].

\bigskip
\noindent\textbf{References:}

[BCM26] M. Buci\'{c}, K. Chen, and J. Ma, On a maximal anti-Ramsey conjeture of Burr, Erd\H{o}s, Graham, and S\'{o}s. arXiv:2603.18952 (2026).
 
 [BEGS89] Burr, S. A. and Erd\H{o}s, P. and Graham, R. L. and S\'os, V.
T., Maximal anti-{R}amsey graphs and the strong chromatic number. J. Graph Theory (1989), 263--282.

\bigskip
\noindent\small{Source: \url{https://www.erdosproblems.com/809}}
\end{document}
